\input{pdfinfo-en}

\documentclass{article}

% Preamble
\usepackage[english]{babel}
\input{preamble/preamble}

\DeclareMathOperator{\tg}{tg}


% Title
\title{\texorpdfstring{Selected problems of entrance examinations\newline Mekhmat MSU in 1986}
                      {Selected problems of entrance examinations, Mekhmat MSU in 1986}}
\author{}
\date{}


\begin{document}

\maketitle

\vspace{-0.4cm}
The level of difficulty of the problems is marked with asterisks; in parentheses - notes by the compiler of the list.
\vspace{0.5cm}


\problem  % Problem 1
Find all rational solutions to the equation
\begin{equation*}
  y = \sqrt{x^2 + \frac{11}{3}}
\end{equation*}

(A simplified analogue of this problem is in the book:
D. O. Shklyarsky, N. N. Chentsov, I. M. Yaglom
``Selected Problems and Theorems of Elementary Mathematics. Arithmetic and Algebra''
Problem No. 273.)


\problem  % Problem 2
Using only a straightedge, construct a line paraller to two given parallel lines, passing through a given point.

(The problem is a [viceroid [?]], O. Shklyarsky, N. N. Chentsov, I. M. Yaglom
``Selected Problems and Theorems of Planimetry'' problem No. 65b.
The solution of problem 65a in the book is used in as an auxiliary construction [?].)


\problem[{**}]  % Problem 3
Compare
$\mspace{1.5mu}
 \dfrac{\log_3\paren{4} \log_3\paren{6} \log_3\paren{8} \cdots \log_3\paren{80}}
       {2 \log_3\paren{5} \log_3\paren{7} \log_3\paren{9} \cdots \log_3\paren{79}}
\mspace{1.5mu}$
and $1$.


\problem  % Problem 4
Solve the equation
\begin{equation*}
  \paren{x - y}^2 + \paren{y - 2 \sqrt{x} + 2}^2 = \frac{1}{2}
\end{equation*}


\problem[{**}]  % Problem 5
The base of the triangular pyramid $DABC$ is the equilateral triangle $ABC$.
Prove that if $\widehat{DAB} \cong \widehat{DBC} \cong \widehat{DCA}$,
then the pyramid is regular.

(The problem is in the problem book of ``Kvant'' (10, 1982), Problem No. 770, with an asterisk, i.e., of greater difficulty.
One month is given to solve this section of the book.)


\problem[{*}]  % Problem 6
Circle $O_1$ is drawn through three non-adjacent vertices of a cube.
Circle $O_2$ is inscribed in a face of the same cube that contains exactly one of these vertices.
Find the distance between the circles $O_1$ and $O_2$.

(An analogous problem was proposed at the final round of the 1985 All-Soviet-Union Mathematical Olympiad (for 10th graders). See No. 1 1985, N 11).

% Clearer formulation: \newline
% Let $ABCD$ be a base of a cube.
% Let $A'B'C'D'$ be the opposite base of that cube,
% with $A'$ adjacent to $A$,
% $B'$ adjacent to $B$,
% $C'$ adjacent to $C$,
% and $D'$ adjacent to $D$.
% Detetrmine the distance between the circle inscribed in $ABCD$
% and the circle circumscribed about $A$, $B'$ and $C$.


\problem  % Problem 7
Let $ABC$ be a triangle.
Construct equilateral triangles $ABC'$, $BCA'$, $CAB'$
externally from $ABC$, using its sides as bases.
Prove that the lines $\overline{AA'}$, $\overline{BB'}$, $\overline{CC'}$ are concurrent.

(This is the main lemma, essentially equivalent to the solution of problem No. 71b*
from the book by D. O. Shklyarsky, N. N. Chentsov, I. M. Yaglom ``Selected Problems and Theorems of Plane Geometry'',
a problem with an asterisk, i.e. of greater difficulty.) \newline
(V. V. Prosimov, Zatan N menuque, 9. 16)


\problem  % Problem 8
Solve the following equation in rational numbers
\begin{equation*}
  \sqrt{2 \sqrt{3} - 3} = \sqrt{x \sqrt{3}} - \sqrt{y \sqrt{3}}
\end{equation*}


\problem  % Problem 9
The median, angle bisector, and height from one vertex of a triangle
divide the corresponding angle into 4 equal parts.
Determine the angles of the triangle.

(The problem is in the book by D. O. Shklyarsky, N. N. Chentsov, I. M. Yaglom ``Selected Problems and theorems in Plane Geometry'', problem No. 113b).


\problem  % Problem 10
Find all triangles that have an inscribed circle of radius $1$
and whose heights are integer.


\problem  % Problem 11
Solve the equation
$x^2 - 5 = \sqrt{x + 5}$.

(A similar problem is in the book:
D. O. Shklyarsky, N. N. Chentsov, I. M. Yaglom
``Selected problems and theorems of elementary mathematics. Arithmetic and Algebra.''
Problem No. 239.)


\problem  % Problem 12
Prove that if
$\dfrac{1}{\mspace{1mu} a + b + c \mspace{1mu}} = \dfrac{1}{a} + \dfrac{1}{b} + \dfrac{1}{c}$,
then, for any $n$, it holds that
$\dfrac{1}{\mspace{1mu} a^n + b^n + c^n \mspace{1mu}} = \dfrac{1}{a^n} + \dfrac{1}{b^n} + \dfrac{1}{c^n}$.

(The main auxiliary statement to this problem is in the book:
V. B. Lidsky, D. V. Ovsyannikov, A. N. Tulaikov, M. I. Shabunin
``Problems in Elementary Mathematics'', No. 252.
Within each subsection, the problems are ordered by their difficulty (increasing from the preface to the end).
The numbers in in the corresponding subsection: 234 - 260.
At the written exam at MITP, where this problem was proposed, five hours were given to solve four problems.)


\problem[{*}]  % Problem 13
Solve the equation
$x^4 + x^3 - 3 a x^2 - 2 a x + 2 a^2 = 0$,
for $x > 0$, $a > 1$.


\problem  % Problem 14
Construct a tangent to the curve $y = \sqrt[3]{x}$
using only straightedge and compass. \newline
?[? Through a given point? or just one tangent will do, whichever it is?]

(The statement can be interpreted in two ways:
the first is that the problem is to be solved with the graph of the curve given;
in this case, the solution is simple.
The second is that the graph of the curve is not given,
only coordinate axes with unit segments on them are given.
Then, solving the problem for a point with abscissa $2$ is a equivalent to the problem of the doubling of a cube,
which is unsolvable with straightedge and compass.
Examiners used both interpretations of the condition.)


\problem  % Problem 15
%\problem[{\footnote{%
%  \label{merged_problems}%
%  Note: during the appeal it became known that the examiners considered these problems as one single problem.%
%  (The problems were offered on the same picture).}}]
A triangle $ABC$ (with $\overline{AB} \cong \overline{BC}$)
is inscribed in a circle so that the bisectors of the external angles at $B$ and $C$ intersect on the circle:
find the radius of the circle.

(The conditions are impossible.)


\problem  % Problem 16
% \problem[{\footref{merged_problems}}]
In a triangle $ABC$, the bisectors of the angles at $B$ and $C$ are drawn.
Does there exist a circle passing through the intersection point of these bisectors,
their bases [the vertices $B$ and $C$] and the vertex $A$?
If so, determine the type of triangle.


\problem  % Problem 17
% \problem[{\footref{merged_problems}}]
A circle is drawn through the bases of two angle bisectors
$\overline{BD}$ and $\overline{CK}$ of triangle $ABC$,
their point of intersection $P$,
and the third vertex $A$ of triangle $ABC$.
$\overline{DK} = 1$.
Determine $\overline{AD}$ and $\overline{AK}$.

(Insufficient data.
Solve the problem with the additional value $\widehat{ABC} = \alpha$.)



%TODO


\end{document}
